%%--------------- 
\begin{ex}%[GVBS: Mai Trung Hiếu]%[2H5V2-7]
Trong không gian $Oxyz$, cho đường thẳng $\Delta \colon \dfrac{x}{2}=\dfrac{y}{2}=\dfrac{z}{1}$ và mặt phẳng $(P): x+2 y-2 z=0$.
Gọi $(Q)$ là mặt phẳng chứa $\Delta$ sao cho góc giữa hai mặt phẳng $(P)$ và $(Q)$ là nhỏ nhất. Phương trình mặt phẳng $(Q)$ là
\choice
{$x-2 y+z=0$}
{$x+22 y+10 z=0$}
{$x-2 y-z=0$}
{\True $x+10 y-22 z=0$}
\loigiai{
\begin{center}
\begin{tikzpicture}[
	>=stealth, line join=round, line cap=round, 
	scale=.8,
	font=\footnotesize,
	declare function={r=3; a=75;}]
% ve diem
	\path 
		(0,0) coordinate (a)
		(2,4) coordinate (b)
		(10,4) coordinate (c)
		(8,0) coordinate (d)
		(4,6) coordinate (e)
		(6,10) coordinate (f)
		(6.5,6) coordinate (B)
		(8.5,1) coordinate (A)
		(9.5,3) coordinate (K)
		(6.5,3) coordinate (H)
		(intersection of b--c and e--d) coordinate (g)
		;
	\draw[thick] 
		(a)--(d)--(c)--(f)--(e)--(d)
		(a)--(b)--(g)
		(K)--(B)
		;
	\draw[thick, shorten >=-1cm, shorten <=-1cm]
		(B)--(A) node[pos=-.2, above right] {$\Delta$}
		;
	\draw[dashed]
		(c)--(g)
		(K)--(H)--(A)
		(B)--(H)
		;
% danh dau goc
	\foreach \a/\b/\c in {c/K/b, K/H/B}	{
		\draw pic [draw=black,angle radius = 6] {right angle = \a--\b--\c};
	}
%	\draw pic ["$1$", draw, angle eccentricity=1.3, angle radius=20] {angle = B--A--C};
% danh dau diem
	\foreach \x/\g in {A/0, B/45, K/0, H/135}{
		\draw[fill=black] (\x) circle (1.5pt)+(\g:.4) node{$\x$};}
	\node[above right] at (a) {$P$};
	\node[below] at (f) {$Q$};
	\node[right] at (9,2) {$d$};
\end{tikzpicture}
\end{center}
Đường thẳng $\Delta \colon \dfrac{x}{2}=\dfrac{y}{2}=\dfrac{z}{1}$ được viết lại dưới dạng tham số $\Delta \colon \heva{&x=2 t \\ &y=2 t \\ &z=t.}$
\\
Xét hệ phương trình
$$
	\heva{&x=2 t \\ &y=2 t \\ &z=t \\ &x+2 y-2 z=0}
	\Leftrightarrow
	\heva{&t=0 \\ &x=0 \\ &y=0 \\ &z=0.}
$$
Do đó $\Delta$ cắt $(P)$ tại điểm $A(0 ; 0 ; 0) \equiv O$.
\\
Lại có $\Delta$ và $(P)$ không vuông góc nhau nên ta đi chứng minh góc nhỏ nhất giữa $(P)$ và $(Q)$ là góc giữa $\Delta$ và $(P)$. Thật vậy trên $\Delta$ lấy $B$ khác $A$, kẻ $BH$ vuông góc với $(P)$ tại $H$ và $B K$ vuông góc $d$ tại $K$ ($d$ là giao tuyến của $(P)$ và $(Q)$) tại $K$. Khi đó góc giữa $(Q)$ và $(P)$ là góc $BKH$.
Ta thấy
$$
	HA \geq HK 
	\Rightarrow 
	\tan \widehat{BKH}=\dfrac{BH}{HK}
		\geq
			\dfrac{BH}{HA}=\tan \widehat{BAH}.
$$
Mà $\widehat{BKH}$, $\widehat{BAH}<90^{\circ}$ nên
$$
	\widehat{BKH} \geq \widehat{BAH}=(\Delta,(P))=\arcsin \dfrac{4}{9}.
$$
Đẳng thức xảy ra $\Leftrightarrow K \equiv A \Leftrightarrow \Delta \perp d$.
\\
Do đó, góc giữa hai mặt phẳng $(P)$ và $(Q)$ là nhỏ nhất khi và chỉ khi $(Q)$ chứa $(\Delta)$ và cắt $(P)$ theo một giao tuyến vuông góc $\Delta$.
\\
*Viết phương trình của $(Q)$
\\
Đường thẳng $\Delta$ có vec-tơ chỉ phương $\vec{u_1}=(2 ; 2 ; 1)$, $(P)$ có vec-tơ pháp tuyến $\vec{n_1}=(1 ; 2 ;-2)$ nên $d$ có vec-tơ chỉ phương $\vec{u_2}=\left[\vec{u_1}, \vec{n_1}\right]=(-6 ; 5 ; 2)$.
\\
$(Q)$ chứa $\Delta$ và $d$ nên nhận $\vec{n_2}=\left[\vec{u_2} ; \vec{u_1}\right]=(1 ; 10 ;-22)$ làm vec-tơ pháp tuyến.
\\
Vậy mặt phẳng $(Q)$ đi qua $A(0 ; 0 ; 0)$ và nhận $\vec{n_2}=(1 ; 10 ;-22)$ làm vec-tơ pháp tuyến nên có phương trình $x+10 y-22 z=0$.
}
\end{ex}

%%--------------- 
\begin{ex}%[GVBS: Mai Trung Hiếu]%[2H5V2-7]
Trong không gian $Oxyz$, cho đường thẳng $d: \dfrac{x}{-1}=\dfrac{y+1}{2}=\dfrac{z-2}{1}$ và mặt phẳng $P \colon 2 x-y-2 z-2=0$.
$(Q)$ là mặt phẳng chứa $d$ và tạo với $(P)$ một góc nhỏ nhất. Gọi $\vec{n}_Q=a ; b ; 1$ là một vec-tơ pháp tuyến của $(Q)$. Đẳng thức nào đúng?
\choice
{$a-b=-1$}
{\True $a+b=-2$}
{$a-b=1$}
{$a+b=0$}
\loigiai{
Đường thẳng $d$ đi qua điểm $M (0 ;-1 ; 2)$ và có vec-tơ chỉ phương $\vec{u_d}=(-1 ; 2 ; 1)$.
\\
Theo giả thiết, $d \subset (Q)$ và $\vec{n}_Q=(a ; b ; 1)$ là một vec-tơ pháp tuyến của $Q$ nên ta có
$$
	\vec{u_d} \cdot \vec{n}_Q=0 \Leftrightarrow-a+2 b+1=0 \Leftrightarrow a=2 b+1 \quad (1).
$$
Mặt phẳng $(P)$ có vec-tơ pháp tuyến $\vec{n}_P=(2 ;-1 ;-2)$.
\\
Ta có
$$
	\cos \left((P), (Q)\right)=\left|\cos \left(\vec{n}_P, \vec{n}_Q\right)\right|=\dfrac{|2 a-b-2|}{3 \cdot \sqrt{a^2+b^2+1}} \quad (2).
$$
Thế (1) vào (2), ta được $\cos \left((P), (Q)\right)=\dfrac{|b|}{\sqrt{5 b^2+4 b+2}}$.
\\
Khi góc giữa $(P)$ và $(Q)$ nhỏ nhất thì $\cos \left((P), (Q)\right)$ đạt giá trị lớn nhất.
\\
Xét hàm số $f(b)=\dfrac{b}{\sqrt{5 b^2+4 b+2}}$, có
$$
	f'(b)=\dfrac{b+1}{\sqrt{5 b^2+4 b+2}}=0 \Leftrightarrow b=-1.
$$
Bảng biến thiên
\begin{center}
\begin{tikzpicture}
	\tkzTabInit[lgt=2,espcl=4,deltacl=0.6]
	{$b$/1, $f'(b)$/1, $f(b)$/2}
	{$-\infty$, $-1$, $+\infty$}
	\tkzTabLine{, -, z, +}
	\tkzTabVar{+/$-\dfrac{1}{\sqrt{5}}$, -/$-\dfrac{1}{\sqrt{3}}$, +/$\dfrac{1}{\sqrt{5}}$}
\end{tikzpicture}
\end{center}
Từ đó suy ra với hàm số $g(b)=\dfrac{|b|}{\sqrt{5 b^2+4 b+2}}$ có
$$
	\max\limits_{\mathbb{R}} g(b)=g(-1)=-\dfrac{1}{\sqrt{3}}.
$$ 
Khi $b=-1 \Rightarrow a=-1$.
\\
Vậy, $a+b=-2$.
}
\end{ex}
\Closesolutionfile{ans}
\indapan{6}{ans/ans-0-B15-KQ}